\section{Equipos arbitrales}


\subsection{Presentación del caso}
Para dirigir 3 partidos de fútbol (P1, P2 y P3) se dispone de 4 equipos arbitrales (A1, A2, A3 y A4). La distancia (en centenares de kilómetros) que debería recorrer cada equipo arbitral para dirigir cada partido es la siguiente:

\begin{center}
\begin{tabular}{c|ccc}
 & P1 & P2 & P3 \\
\hline
A1 & 5 & 7 & 7 \\
A2 & 7 & 5 & 8 \\
A3 & 6 & 4 & 8 \\
A4 & 7 & 6 & 9 \\
\end{tabular}
\end{center}

Se trata de determinar qué equipo debe arbitrar cada partido de modo que la distancia total recorrida en la jornada de liga sea mínima.

\subsection{Formulación explícita}

\paragraph{Variables}\mbox{}\\
\begin{tabular}{l c l}
    $x_{ij}$ & : & variable binaria que toma valor 1 si el equipo arbitral $i$ dirige el partido $j$,\\
    & & 0 en caso contrario.\\
\end{tabular}

\paragraph{Restricciones}\mbox{}\\
Cada partido debe ser dirigido por exactamente un equipo arbitral:
\begin{align}
x_{11}+x_{21}+x_{31}+x_{41} &= 1 \\
x_{12}+x_{22}+x_{32}+x_{42} &= 1 \\
x_{13}+x_{23}+x_{33}+x_{43} &= 1
\end{align}
\\
Cada equipo arbitral puede dirigir como máximo un partido:
\begin{align}
x_{11}+x_{12}+x_{13} &\leq 1 \\
x_{21}+x_{22}+x_{23} &\leq 1 \\
x_{31}+x_{32}+x_{33} &\leq 1 \\
x_{41}+x_{42}+x_{43} &\leq 1
\end{align}
\\

\paragraph{Función objetivo}\mbox{}\\
Minimizar la distancia total recorrida (en centenares de km):
\begin{equation}
\mbox{min. } z = 5x_{11}+7x_{12}+7x_{13}+7x_{21}+5x_{22}+8x_{23}+6x_{31}+4x_{32}+8x_{33}+7x_{41}+6x_{42}+9x_{43}
\end{equation}

\paragraph{Modelo completo}\mbox{}\\
\begin{align}
\mbox{min.} \quad 
& z = 5x_{11}+7x_{12}+7x_{13}+7x_{21}+5x_{22}+8x_{23}+6x_{31}+4x_{32}+8x_{33}+7x_{41}+6x_{42}+9x_{43} \notag \\
\mbox{s. a.:} \quad 
& x_{11}+x_{21}+x_{31}+x_{41} = 1 \notag \\
& x_{12}+x_{22}+x_{32}+x_{42} = 1 \notag \\
& x_{13}+x_{23}+x_{33}+x_{43} = 1 \notag \\
& x_{11}+x_{12}+x_{13} \leq 1 \notag \\
& x_{21}+x_{22}+x_{23} \leq 1 \notag \\
& x_{31}+x_{32}+x_{33} \leq 1 \notag \\
& x_{41}+x_{42}+x_{43} \leq 1 \notag \\
& x_{ij} \in \{0,1\} \quad \forall i=1,\dots,4;\; j=1,2,3 \notag
\end{align}


\subsection{Formulación generalizada}

\paragraph{Conjuntos}\mbox{}\\
\begin{tabular}{l c l}
    $\mathcal{A}$ & : & conjunto de equipos arbitrales ($i \in \mathcal{A}$)\\
    $\mathcal{P}$ & : & conjunto de partidos ($j \in \mathcal{P}$)\\
\end{tabular}
\\

\paragraph{Parámetros}\mbox{}\\
\begin{tabular}{l c l}
    $D_{ij}$ & : & distancia (en centenares de km) que debe recorrer el equipo $i$ para dirigir el partido $j$ \\
\end{tabular}
\\

\paragraph{Variables}\mbox{}\\
\begin{tabular}{l c l}
    $x_{ij}$ & : & variable binaria que vale 1 si el equipo $i$ arbitra el partido $j$, 0 en caso contrario \\
\end{tabular}
\\

\paragraph{Restricciones}\mbox{}\\
Cada partido debe tener un único equipo asignado:
\begin{equation}
    \sum_{i \in \mathcal{A}} x_{ij} = 1 \quad \forall j \in \mathcal{P}
\end{equation}
\\
Cada equipo arbitral puede dirigir como máximo un partido:
\begin{equation}
    \sum_{j \in \mathcal{P}} x_{ij} \leq 1 \quad \forall i \in \mathcal{A}
\end{equation}
\\

\paragraph{Función objetivo}\mbox{}\\
\begin{equation}
\mbox{min. } z = \sum_{i \in \mathcal{A}} \sum_{j \in \mathcal{P}} D_{ij}\, x_{ij}
\end{equation}

\paragraph{Modelo completo}\mbox{}\\
\begin{align}
\mbox{min.} \quad & z = \sum_{i \in \mathcal{A}} \sum_{j \in \mathcal{P}} D_{ij}\, x_{ij} \notag \\
\mbox{s. a.:} \quad & \sum_{i \in \mathcal{A}} x_{ij} = 1 \quad \forall j \in \mathcal{P} \notag \\
& \sum_{j \in \mathcal{P}} x_{ij} \leq 1 \quad \forall i \in \mathcal{A} \notag \\
& x_{ij} \in \{0,1\} \quad \forall i \in \mathcal{A},\; j \in \mathcal{P} \notag
\end{align}
